\section{Interpretation}
\label{sec:interpretation}

This section identifies which results are novel relative to the literature cited in
Section~\ref{sec:intro}, why each is of interest, what confound remains, and what would
establish it. Results are ordered by how far they depart from what the existing literature
would predict.

\subsection{Agreement calibration reverses sign across task families}

\textbf{Finding.} $\dvote$ is $+0.15$ to $+0.19$ on GPQA and $+0.05$ to $+0.09$ on
TruthfulQA, but $-0.16$ to $-0.19$ on MMLU-Pro. The same three-atom agreement signal, under
the same aggregation rule and the same models, is substantially overconfident on one task
family and substantially underconfident on another.

\textbf{Why this is novel.} Neither side of the existing debate predicts a sign reversal.
\citet{yang2024collaborative} report that multi-agent deliberation improves confidence
assessment, which predicts a negative delta. \citet{wynn2025debate} document
consensus on wrong answers, which predicts a positive delta. Both are correct in this
sweep, on different benchmarks, with everything else held fixed. The disagreement in the
literature is therefore not necessarily a disagreement about mechanism; it may be a
disagreement about the evaluation suite.

\textbf{Why it is interesting.} It means agreement calibration is not a property of a
coordination protocol at all. It is a property of the interaction between the protocol and
the answer-space geometry. The most plausible mechanism is option-set cardinality: GPQA
Diamond has four options, so unanimous agreement among three agents is common and cheap,
while MMLU-Pro has up to ten, so unanimity is rare and correspondingly more informative.
Under that account the vote fraction is not a calibrated probability in either case; it is
a statistic whose relationship to correctness depends on how many ways there are to be
wrong.

\textbf{What would establish it.} The cardinality account is testable and the sweep does not
test it. Holding the benchmark fixed and varying the number of presented distractors would
separate option cardinality from the many other differences between GPQA and MMLU-Pro,
including domain, difficulty and question construction. A second, cheaper test is to fit a
per-benchmark recalibration map for the vote atoms and check whether the fitted maps differ
in the way cardinality predicts.

\subsection{The coordination calibration penalty shrinks with capacity}

\textbf{Finding.} Vote calibration error falls faster with parameter count than per-agent
calibration error does. The exponent contrasts are $+0.077$ $[0.031, 0.127]$, $+0.495$
$[0.424, 0.567]$ and $+0.087$ $[0.008, 0.153]$ on the three multiple-choice benchmarks,
all excluding zero. At the two largest capacities $\dvote$ turns negative in several
strata.

\textbf{Why this is novel.} The sweep was designed to test the opposite hypothesis, that
coordination penalties compound with capability. \citet{kim2025scaling} frame agent-system
scaling in terms of error amplification and redundancy, quantities that do not obviously
improve with component quality. The direction found here is that better components make
coordination a calibration benefit rather than a cost, and that the crossover happens
within the range of models a practitioner would actually deploy.

\textbf{Why it is interesting.} If it holds, the design advice inverts with scale. At small
capacity, exposing agreement as confidence is harmful and should be suppressed. At large
capacity it becomes a useful signal that the individual models do not provide. That is a
qualitatively different recommendation from a monotone one, and it means results obtained
on small open models do not extrapolate.

\textbf{What confound remains.} Vote-atom coarseness is the leading alternative
explanation. With three agents the signal has three atoms, and a change in how probability
mass distributes over three atoms can move the exact-atom calibration error without any
change in the honesty of agreement. The final-producer signal, which has continuous support
and does not suffer from this, shows no comparable capacity effect, which is what the
coarseness explanation predicts.

\textbf{What would establish it.} Repeating the measurement at five and seven agents, which
gives six and eight atoms, would separate the two accounts directly: a genuine calibration
improvement should persist as the support refines, whereas a coarseness artefact should
attenuate. The frozen manifests for those runs already exist.

\subsection{Agreement calibration and producer calibration are different quantities}

\textbf{Finding.} $\dvote$ spans a range of roughly 0.35 across benchmarks and topologies.
$\dfp$ spans $+0.004$ to $+0.020$ across every combination where it is defined.

\textbf{Why this is novel.} The multi-agent calibration literature reports a single
system-level calibration number. Separating the honesty of the agreement statistic from the
calibration of the decisive model is, as far as the cited work goes, not done, and the
two answers differ by an order of magnitude on the same data.

\textbf{Why it is interesting.} It localises the problem. Coordination is not corrupting the
underlying models: the model that produces the answer remains almost as well calibrated as
it was alone. What is miscalibrated is the conversion of agreement into a number. That is a
much more tractable failure, because it is a property of the aggregation rule rather than of
the models, and aggregation rules are cheap to change.

The practical form of the claim is direct. A system that displays how many agents agreed
can be wrong in either direction by 0.15 or more depending on the task. A system that
displays the decisive model's own probability incurs a penalty small enough to be removed by
post-hoc scaling. The second is strictly preferable and costs nothing extra to compute.

\textbf{What would establish it further.} Fitting a held-out confidence combiner over the
vote fraction, the producer probability, the agreement level and the verbalized signal would
test whether the aggregation rule is recoverable, and would give an upper bound on how much
calibration a well-chosen rule can recover.

\subsection{Reasoning improves accuracy and moves confidence past calibration}

\textbf{Finding.} Realized reasoning tokens raise accuracy, with dose slopes of $+0.0106$ to
$+0.0175$, and reduce signed overconfidence, with slopes of $-0.0058$ to $-0.0140$, while
simultaneously raising per-agent calibration error, with slopes of $+0.0034$ to $+0.0126$.

\textbf{Why this is novel.} \citet{zhu2025testtime} establish that test-time compute
improves agent performance, and \citet{li2026agentbench} benchmark that scaling. Neither
reports a calibration cost. The result here is that the accuracy gain is real and the
overconfidence reduction is real, but they are not the same thing as improved calibration,
and on this sweep the raw probability signal gets worse as the model thinks longer.

\textbf{Why it is interesting.} The natural reading is that extended reasoning shifts the
confidence distribution downward past the calibration point rather than sharpening it onto
that point, converting overconfidence into underconfidence. If that is what is happening,
the fix is a monotone recalibration rather than any change to the reasoning mechanism, and
post-hoc temperature scaling should recover most of the loss.

\textbf{What would establish it.} The reliability diagrams and the Cox slopes already
computed per cell bear on this directly: an underconfidence account predicts Cox slopes
moving in a specific direction with dose. Comparing the raw and temperature-scaled
calibration errors as a function of dose would settle whether the loss is a monotone
distortion. A confound to remove first is that the thinking-enabled and thinking-disabled
conditions use different sampling temperatures, 0.6 and 0.7, so part of the observed
confidence shift may be a sampling effect rather than a reasoning effect.

\subsection{Confident-wrong consensus attenuates with capacity}

\textbf{Finding.} Debate raises unanimous incorrect consensus relative to independent voting
by $+0.0347$ at 0.6B and 1.7B, $+0.0172$ at 4B and 8B, and $+0.0164$ at 14B and 32B.

\textbf{Why this refines the literature.} \citet{wynn2025debate} identify the failure mode.
This sweep confirms it at every capacity and adds that its magnitude more than halves across
a fifty-fold range of parameters. The failure is therefore real but is a small-model failure
in the main, which changes how seriously it should be weighted when choosing a protocol for
a large model.

\textbf{Why it is interesting.} It coexists with debate having the best average accuracy and
the lowest error amplification of the three multi-agent topologies. Debate reduces the total
error rate while increasing the rate of a specific high-cost failure in which the system is
wrong and maximally confident. That trade is acceptable when only the answer is exposed and
unacceptable when agreement is exposed as confidence.

\textbf{What confound remains.} Decentralized runs two rounds and independent runs one, so
the contrast includes an extra round of generation. A one-round decentralized ablation
separates communication from compute and is the single cheapest experiment that would
sharpen this result.

\subsection{The obvious system confidence signal is the worst one}

\textbf{Finding.} Mean agreeing option-logprob confidence wins 0 of 18 strata in the signal
tournament. Mean per-agent confidence over all agents wins 13, and mean agreeing verbalized
confidence wins 5.

\textbf{Why it is interesting.} Averaging the probabilities of the agents that agreed is the
construction a practitioner would reach for first, and it is beaten in every stratum by not
conditioning on agreement at all. Conditioning selects agents that were confident in the
winning answer, which raises the mean without adding information, and the selection is
strongest exactly when the team is unanimous.

That verbalized confidence beats the logprob signal in five strata is also worth noting,
given that verbalized confidence is usually treated as the weaker signal and is far cheaper
to obtain since it requires no logprob access.

\textbf{What would establish it.} The tournament is a comparison of fixed signals. A learned
combiner would establish whether the ordering reflects the information content of the
signals or only the particular functional forms tested.

\subsection{Matched-compute multi-agent systems lose on knowledge tasks and win on mathematics}

\textbf{Finding.} Three 8B agents against one 32B agent: better in 3 of 12 contrasts, worse
in 7. Mean differences are $-0.121$ on GPQA, $-0.117$ on TruthfulQA, $-0.034$ on MMLU-Pro
and $+0.059$ on MATH.

\textbf{Why this matters.} \citet{smit2024debate} argue that debate is not reliably superior
at matched cost. This result agrees and adds the task-family split: the exception is
mathematics, and it is a large exception.

\textbf{Why it is interesting.} The split has a plausible mechanism that the rest of the
sweep supports. Aggregation helps when solutions are long, independently derivable and
checkable against each other, which describes mathematical derivation. It does not help when
the task is recall of a specific fact, where three smaller models are three draws from the
same insufficient knowledge rather than three independent attempts. Every other topology
result in the sweep points the same way: independent voting gains $+0.076$ on mathematics,
its largest effect anywhere, and centralized orchestration loses $0.071$ on the same
benchmark, its largest loss anywhere.

\textbf{What confound remains.} Some 32B long-reasoning baselines are context-capped, which
depresses the single-agent side of the comparison in precisely the configurations where it
should be strongest, so the advantage of the larger model is likely understated. A three-4B
against one-14B comparison, where neither side is capped, would be cleaner.

\subsection{Peer rationales help only in a specific regime, but cost calibration everywhere}

\textbf{Finding.} The controlled plus-CoT accuracy effect is $+0.045$ with 16 of 16 strata
positive in the mid-capacity low-reasoning band, and $+0.006$ with 11 positive against 12
negative in the large-capacity high-reasoning band. The calibration cost does not vary in
this way: vote calibration error rises in 94 of 119 strata regardless.

\textbf{Why it is interesting.} The accuracy benefit and the calibration cost have different
regime dependence, which means they are not the same mechanism. The benefit behaves like an
information transfer that saturates once the receiving model can generate equivalent content
itself. The cost behaves like a convergence effect that operates whenever agents see each
other's reasoning, whether or not that reasoning was useful.

The item-level evidence supports this reading: exposing chain of thought tightens the spread
of confidence across agents in five of six strata without improving individual calibration.
Agents converge; the convergence is not evidence.

\textbf{What would establish it.} Comparing compressed rationales against raw chains of
thought would separate the useful information from the social-proof effect. If compressed
rationales retain the accuracy benefit and lose the calibration cost, the two mechanisms are
confirmed as separable and the design recommendation is straightforward.

\subsection{A negative result worth recording}

The voting-lift decomposition finds no established error-correlation penalty: the within-cell
penalty interval is below zero in 0 of 24 strata. This is a measurement limitation rather
than evidence of absence. Option letters are shuffled independently per question and seed,
so two agents selecting the same wrong content in different cells appear as different
letters, and cross-seed wrong-answer collision rates are not identifiable.

This is worth recording because the item-level confident-wrong analysis does find a
correlation effect using a different construction. The two results are not in conflict; the
decomposition simply cannot be computed from what was logged. Recording canonical answer
identities alongside the positional letters would make it computable and costs nothing at
run time.

\subsection{Summary of the design implications}

The results support four statements that are stronger than the individual effects because
several independent lines converge on each.

Capacity is the dominant lever, worth roughly four times any other axis in the sweep, and
prompt scaffolding is the most cost-effective, being the only intervention that improves
both accuracy and agreement calibration on every benchmark at negligible cost.

Reasoning budget should be set at b2048 or b8192 and not higher. The unlimited rung is zero
or negative against b8192 on every benchmark, and the b512 rung is worse than no thinking at
all on mathematics.

Orchestration should not be the default. Centralized orchestration does not beat debate on
accuracy or voting on confidence while costing as much as debate, and no multi-agent topology
recovers its turn cost under coordination efficiency.

If a system exposes confidence, it should expose a recalibrated producer probability rather
than an agreement fraction. This is the clearest and best-supported design conclusion in the
report, and it follows from the order-of-magnitude difference between the two deltas rather
than from any single effect estimate.
